\documentclass[11pt]{article}
\usepackage{amsmath,amssymb,bm}
\usepackage[a4paper,margin=1in]{geometry}
\usepackage{mathtools}

\title{Transfer-function derivation with eddy-viscosity closure}
\author{}
\date{}

\begin{document}
\maketitle

\section{Problem setting}

We use the same coordinates as in the original note:
\[
x:\text{streamwise}, \qquad y:\text{spanwise}, \qquad z:\text{wall-normal},
\]
and impose the streamwise-constant assumption
\[
\partial_x(\cdot)=0.
\]
The wall-normal derivative is denoted by
\[
D := \frac{\mathrm d}{\mathrm d z}.
\]

The starting point is the Craik--Leibovich (CL) equation in non-dimensional form,
\begin{equation}
\frac{\partial \bm u}{\partial t}
+(\bm u\cdot\nabla)\bm u
=-\nabla p
+\frac{1}{La}\,\bm u_s\times \bm\omega
+\frac{1}{Re}\nabla^2\bm u
+\bm f,
\qquad
\nabla\cdot\bm u=0,
\label{eq:cl-original}
\end{equation}
where
\[
\bm\omega=\nabla\times\bm u,
\qquad
\bm u_s=u_s(z)\,\bm e_x.
\]

In the original article, the nonlinear fluctuation term is absorbed into the external forcing. Here we do \emph{not} do that. Instead, we model the nonlinear term with an eddy viscosity.

\section{Eddy-viscosity closure of the nonlinear term}

Let the velocity be decomposed as
\[
\bm u = U(z)\,\bm e_x + \bm u',
\qquad
\bm u' = (u,v,w),
\]
where \(U(z)\) is the base streamwise velocity. The exact perturbation equation contains the nonlinear term
\[
\bm{\mathcal N}
:=
-(\bm u'\cdot\nabla)\bm u'
\quad
\text{(or equivalently the divergence of the perturbation Reynolds stress).}
\]

We close it by a Boussinesq-type eddy-viscosity model:
\begin{equation}
\mathcal N_i
\approx
\partial_j\!\left[
\nu_t(z)\left(
\partial_j u_i + \partial_i u_j
\right)
\right].
\label{eq:boussinesq-closure}
\end{equation}
After adding the molecular viscosity, define the total effective viscosity
\begin{equation}
\nu_e(z)
:=
\frac{1}{Re}+\nu_t(z).
\label{eq:nu-effective}
\end{equation}
Then the perturbation equation becomes
\begin{equation}
\frac{\partial \bm u'}{\partial t}
+U\,\partial_x \bm u'
+w\,U'(z)\,\bm e_x
=
-\nabla p'
+\frac{1}{La}\,\bm u_s\times(\nabla\times \bm u')
+\nabla\cdot\left[
\nu_e(z)\left(\nabla \bm u'+\nabla \bm u'^{\,T}\right)
\right]
+\bm f,
\label{eq:linear-eddy-model}
\end{equation}
with
\[
\nabla\cdot \bm u'=0.
\]

\paragraph{Important remark.}
Equation \eqref{eq:linear-eddy-model} is linear \emph{only if} \(\nu_t\) is treated as a prescribed profile, for example \(\nu_t=\nu_t(z)\). This is the usual ``frozen eddy viscosity'' assumption. If instead \(\nu_t\) is itself a nonlinear functional of the instantaneous perturbation field, such as a Smagorinsky model, then the system remains nonlinear and the transfer-function/eigenvalue framework below is no longer exact unless one linearizes \(\nu_t\) about a reference state.

\section{Component form under the streamwise-constant assumption}

Because \(\partial_x=0\), the continuity equation reduces to
\begin{equation}
\partial_y v + D w = 0.
\label{eq:continuity-physical}
\end{equation}

For \(\nu_e=\nu_e(z)\), the modeled viscous operator is
\[
\left[\nabla\cdot\left(
\nu_e(\nabla \bm u'+\nabla \bm u'^{\,T})
\right)\right]_i
=
\nu_e \nabla^2 u_i
+\nu_e' \left(Du_i+\partial_i w\right).
\]
Therefore, the three momentum equations become
\begin{subequations}
\begin{align}
\frac{\partial u}{\partial t} + w U'
&=
\nu_e\left(\partial_{yy}+D^2\right)u
+\nu_e' D u
+f_x,
\label{eq:u-physical}
\\
\frac{\partial v}{\partial t}
&=
-\partial_y p
+\nu_e\left(\partial_{yy}+D^2\right)v
+\nu_e'\left(Dv+\partial_y w\right)
+\frac{1}{La}\,u_s\,\partial_y u
+f_y,
\label{eq:v-physical}
\\
\frac{\partial w}{\partial t}
&=
-D p
+\nu_e\left(\partial_{yy}+D^2\right)w
+2\nu_e' D w
+\frac{1}{La}\,u_s\,D u
+f_z.
\label{eq:w-physical}
\end{align}
\end{subequations}

Compared with the constant-viscosity derivation, the change is clear:
\begin{itemize}
\item the nonlinear term no longer enters as part of the forcing;
\item instead it changes the linear operator through \(\nu_e(z)\), \(\nu_e'(z)\), and later also \(\nu_e''(z)\);
\item the CL coupling \(u \leftrightarrow w\) caused by the vortex-force term is still present.
\end{itemize}

\section{Fourier transform in the spanwise direction}

Introduce the Fourier transform in \(y\):
\[
a(y,z,t)
\; \mapsto \;
\tilde a(k,z,t),
\qquad
\partial_y \mapsto ik,
\qquad
\Delta := D^2-k^2.
\]
Then \eqref{eq:continuity-physical} gives
\begin{equation}
ik\,\tilde v + D\tilde w = 0,
\qquad
\tilde v = -\frac{D\tilde w}{ik}.
\label{eq:continuity-fourier}
\end{equation}

The transformed equations are
\begin{subequations}
\begin{align}
\frac{\partial \tilde u}{\partial t} + U' \tilde w
&=
\underbrace{\left(\nu_e \Delta + \nu_e' D\right)}_{=: \mathcal L_u}\tilde u
+\tilde f_x,
\label{eq:u-fourier}
\\
\frac{\partial \tilde v}{\partial t}
&=
-ik\tilde p
+\nu_e\Delta \tilde v
+\nu_e'\left(D\tilde v + ik\tilde w\right)
+\frac{ik}{La}\,u_s\,\tilde u
+\tilde f_y,
\label{eq:v-fourier}
\\
\frac{\partial \tilde w}{\partial t}
&=
-D\tilde p
+\nu_e\Delta \tilde w
+2\nu_e' D\tilde w
+\frac{1}{La}\,u_s\,D\tilde u
+\tilde f_z.
\label{eq:w-fourier}
\end{align}
\end{subequations}

\section{Modified Orr--Sommerfeld/Squire-type system}

The streamwise equation \eqref{eq:u-fourier} already has the desired form. The nontrivial step is to eliminate \(\tilde p\) and \(\tilde v\) from \eqref{eq:v-fourier}--\eqref{eq:w-fourier}.

\subsection{Pressure elimination}

Apply \(ik\) to \eqref{eq:v-fourier} and \(D\) to \eqref{eq:w-fourier}, then add the two equations. Using
\[
ik\,\frac{\partial \tilde v}{\partial t}
+D\,\frac{\partial \tilde w}{\partial t}
=
\frac{\partial}{\partial t}(ik\tilde v + D\tilde w)
=0
\]
from continuity, we obtain
\begin{equation}
\Delta \tilde p
=
ik\,\mathcal A_v + D\,\mathcal A_w
+ik\tilde f_y + D\tilde f_z
+\frac{ik}{La}\,u_s\,ik\tilde u
+\frac{1}{La}D\!\left(u_s D\tilde u\right),
\label{eq:pressure-poisson}
\end{equation}
where
\begin{equation}
\mathcal A_v := \nu_e\Delta \tilde v + \nu_e'(D\tilde v + ik\tilde w),
\qquad
\mathcal A_w := \nu_e\Delta \tilde w + 2\nu_e' D\tilde w.
\label{eq:Aw-Av}
\end{equation}

Now apply \(\Delta\) to \eqref{eq:w-fourier}:
\begin{equation}
\frac{\partial \Delta \tilde w}{\partial t}
=
-D\Delta \tilde p
+\Delta \mathcal A_w
+\frac{1}{La}\Delta\!\left(u_s D\tilde u\right)
+\Delta \tilde f_z.
\label{eq:Delta-w-before}
\end{equation}
Substituting \eqref{eq:pressure-poisson} into \eqref{eq:Delta-w-before} yields
\begin{equation}
\frac{\partial \Delta \tilde w}{\partial t}
=
-ik\,D\mathcal A_v
-k^2 \mathcal A_w
-ik\,D\tilde f_y
-k^2 \tilde f_z
+\frac{k^2}{La}\,u_s'\tilde u.
\label{eq:Delta-w-intermediate}
\end{equation}

\subsection{Reduction to an operator acting on \(\tilde w\)}

Using \eqref{eq:continuity-fourier},
\[
\tilde v = -\frac{D\tilde w}{ik},
\qquad
\Delta \tilde v = -\frac{D\Delta \tilde w}{ik},
\]
one can reduce the viscous combination in \eqref{eq:Delta-w-intermediate} to
\begin{equation}
-ik\,D\mathcal A_v
-k^2 \mathcal A_w
=
\nu_e \Delta^2 \tilde w
+2\nu_e' \Delta D\tilde w
+\nu_e''(D^2+k^2)\tilde w.
\label{eq:variable-viscous-OS}
\end{equation}
Hence the coupled system becomes
\begin{subequations}
\begin{align}
\frac{\partial \tilde u}{\partial t} + U' \tilde w
&=
\mathcal L_u \tilde u + \tilde f_x,
\label{eq:os-u}
\\
\frac{\partial \Delta \tilde w}{\partial t}
&=
\underbrace{\left[
\nu_e \Delta^2
+2\nu_e' \Delta D
+\nu_e''(D^2+k^2)
\right]}_{=:\mathcal L_w}
\tilde w
-ikD\tilde f_y
-k^2\tilde f_z
+\frac{k^2}{La}\,u_s'\tilde u.
\label{eq:os-w}
\end{align}
\end{subequations}

This is the modified Orr--Sommerfeld/Squire-type formulation with eddy-viscosity closure.

\paragraph{Key difference from the original derivation.}
Originally, the nonlinear term was placed into the forcing and the operators were simply
\[
\mathcal L_u = \frac{1}{Re}\Delta,
\qquad
\mathcal L_w = \frac{1}{Re}\Delta^2.
\]
After eddy-viscosity closure, the forcing no longer contains the modeled nonlinearity; instead the resolvent/eigenvalue operator is replaced by the variable-coefficient operators
\[
\mathcal L_u = \nu_e \Delta + \nu_e' D,
\qquad
\mathcal L_w = \nu_e \Delta^2 + 2\nu_e' \Delta D + \nu_e''(D^2+k^2).
\]

\section{Laplace transform and transfer function}

Take the Laplace transform in time:
\[
\hat a(k,z,s) := \mathcal L_t\{\tilde a(k,z,t)\},
\qquad
\mathcal L_t\left\{\frac{\partial a}{\partial t}\right\}=s\hat a.
\]
Then \eqref{eq:os-u}--\eqref{eq:os-w} become
\begin{subequations}
\begin{align}
\left(sI-\mathcal L_u\right)\hat u + U'\hat w
&=
\hat f_x,
\label{eq:laplace-u}
\\
\left(s\Delta-\mathcal L_w\right)\hat w
&=
-ikD\hat f_y
-k^2\hat f_z
+\frac{k^2}{La}\,u_s' \hat u.
\label{eq:laplace-w}
\end{align}
\end{subequations}

Solving \eqref{eq:laplace-u} for \(\hat u\),
\begin{equation}
\hat u
=
\left(sI-\mathcal L_u\right)^{-1}
\left(
-U'\hat w + \hat f_x
\right).
\label{eq:u-from-w}
\end{equation}
Substitute \eqref{eq:u-from-w} into \eqref{eq:laplace-w}:
\begin{align}
\Bigl[
s\Delta-\mathcal L_w
+\frac{k^2}{La}\,u_s'
\left(sI-\mathcal L_u\right)^{-1}U'
\Bigr]\hat w
&=
-ikD\hat f_y
-k^2\hat f_z
+\frac{k^2}{La}\,u_s'
\left(sI-\mathcal L_u\right)^{-1}\hat f_x.
\label{eq:w-operator-form}
\end{align}

Define
\begin{subequations}
\begin{align}
G_{1,\nu}(s)
&:=
\left(s\Delta-\mathcal L_w\right)^{-1}(-ikD),
\label{eq:G1}
\\
G_{2,\nu}(s)
&:=
\left(s\Delta-\mathcal L_w\right)^{-1}(-k^2),
\label{eq:G2}
\\
G_{3,\nu}(s)
&:=
\frac{k^2}{La}
\left(s\Delta-\mathcal L_w\right)^{-1}
u_s'
\left(sI-\mathcal L_u\right)^{-1},
\label{eq:G3}
\\
\Gamma_\nu(s)
&:=
G_{3,\nu}(s)\,U'.
\label{eq:Gamma}
\end{align}
\end{subequations}
Then
\begin{equation}
\left(I+\Gamma_\nu(s)\right)\hat w
=
G_{1,\nu}(s)\hat f_y
+G_{2,\nu}(s)\hat f_z
+G_{3,\nu}(s)\hat f_x,
\label{eq:w-transfer-pre}
\end{equation}
that is,
\begin{equation}
\hat w
=
\left(I+\Gamma_\nu(s)\right)^{-1}
\Bigl[
G_{1,\nu}(s)\hat f_y
+G_{2,\nu}(s)\hat f_z
+G_{3,\nu}(s)\hat f_x
\Bigr].
\label{eq:w-transfer-final}
\end{equation}

Therefore the total transfer operator is
\begin{equation}
G_\nu(s)
:=
\left(I+\Gamma_\nu(s)\right)^{-1}
\begin{bmatrix}
G_{1,\nu}(s) & G_{2,\nu}(s) & G_{3,\nu}(s)
\end{bmatrix}.
\label{eq:total-transfer}
\end{equation}
Strictly speaking, this is an \emph{operator-valued} transfer function, so the denominator must be written as \(I+\Gamma_\nu(s)\), not merely \(1+\Gamma_\nu(s)\).

\section{Generalized eigenvalue problem}

If we ignore the forcing, the Laplace-domain system gives the generalized eigenvalue problem
\begin{equation}
A_\nu \psi = s\,B\,\psi,
\qquad
\psi=
\begin{bmatrix}
\hat u\\[3pt]
\hat w
\end{bmatrix},
\label{eq:eig-problem}
\end{equation}
with
\begin{equation}
A_\nu
=
\begin{bmatrix}
\mathcal L_u & -U' I\\[3pt]
\dfrac{k^2}{La}\,u_s' I & \mathcal L_w
\end{bmatrix},
\qquad
B=
\begin{bmatrix}
I & 0\\
0 & \Delta
\end{bmatrix}.
\label{eq:eig-matrices}
\end{equation}
This is the direct analogue of your original eigenvalue form, except that
\[
\frac{1}{Re}\Delta
\quad\text{and}\quad
\frac{1}{Re}\Delta^2
\]
must now be replaced by
\[
\mathcal L_u = \nu_e \Delta + \nu_e' D,
\qquad
\mathcal L_w = \nu_e \Delta^2 + 2\nu_e' \Delta D + \nu_e''(D^2+k^2).
\]

\section{Two important simplifications}

\subsection{Case 1: constant eddy viscosity}

If \(\nu_t\) is taken as a constant, then \(\nu_e'=\nu_e''=0\), and the whole derivation collapses to the same algebraic structure as the original note:
\[
\mathcal L_u = \nu_e \Delta,
\qquad
\mathcal L_w = \nu_e \Delta^2.
\]
Equivalently, one just replaces
\[
\frac{1}{Re}
\longrightarrow
\nu_e = \frac{1}{Re}+\nu_t.
\]
So in the constant-\(\nu_t\) case, the later derivation does \emph{not} fundamentally change; only the effective viscous damping changes.

\subsection{Case 2: variable but frozen eddy viscosity \(\nu_t=\nu_t(z)\)}

This is the most useful linear extension. The transfer-function and eigenvalue analysis still works, but:
\begin{itemize}
\item \(\Delta/Re\) is replaced by \(\mathcal L_u\);
\item \(\Delta^2/Re\) is replaced by \(\mathcal L_w\);
\item the resolvent becomes a variable-coefficient operator;
\item additional terms involving \(\nu_e'\) and \(\nu_e''\) modify both damping and modal structure.
\end{itemize}
Therefore, the most amplified spanwise mode and the least stable eigenvalue will generally shift, even if the CL coupling term is unchanged.

\subsection{Case 3: dynamic eddy viscosity \(\nu_t=\nu_t(\bm u')\)}

If \(\nu_t\) depends on the perturbation field itself, then the closure introduces new nonlinearities through \(\delta\nu_t\), and the equations are no longer linear in \((u,v,w)\). In that case:
\begin{itemize}
\item the exact transfer function is no longer defined in the linear sense;
\item the generalized eigenvalue problem \eqref{eq:eig-problem} is no longer exact;
\item one must first linearize \(\nu_t\) around a reference state, e.g.
\[
\nu_t(\bm u') \approx \bar\nu_t(z) + \delta\nu_t,
\]
and retain only the terms linear in the perturbation.
\end{itemize}

\section{Final summary}

If the nonlinear term is closed by an eddy viscosity instead of being absorbed into the forcing, then the subsequent derivation changes in the following way:
\begin{enumerate}
\item the modeled nonlinearity moves from the forcing side to the linear operator side;
\item the streamwise equation changes from
\[
\frac{1}{Re}\Delta \tilde u
\]
to
\[
\mathcal L_u \tilde u
=
\left(\nu_e \Delta + \nu_e' D\right)\tilde u;
\]
\item the wall-normal Orr--Sommerfeld operator changes from
\[
\frac{1}{Re}\Delta^2 \tilde w
\]
to
\[
\mathcal L_w \tilde w
=
\left(\nu_e \Delta^2 + 2\nu_e' \Delta D + \nu_e''(D^2+k^2)\right)\tilde w;
\]
\item the transfer function is still available under the frozen-\(\nu_t\) assumption, but it must be understood as an operator resolvent;
\item the generalized eigenvalue problem remains valid under the frozen-\(\nu_t\) assumption, with \(A\) replaced by \(A_\nu\).
\end{enumerate}

\end{document}
